Badamy rozszerzanie zadanej bijekcji pięciu punktów płaszczyzny $\Q^2$ do automorfizmu wielomianowego
o współczynnikach wymiernych i możliwie najmniejszym stopniu; konstrukcja ogólna daje stopień $(m-1)^2=16$.
Najniższy stopień rozwiązania przyjętego przez niezależny weryfikator wynosi $6$ (ZWERYFIKOWANE): jest to
złożenie czynników trójkątnych stopni $3$ i $2$ z odwzorowaniami afinicznymi, otrzymane przez rozwiązanie
dwóch układów liniowych dla dobrze dobranych punktów pośrednich.
Korzystając z twierdzenia Junga--van der Kulka (cytowanego), pokazujemy, że nad $\Q$ stopień $\le5$
jest niemożliwy; dowód sprowadza się do dokładnych rachunków (jądro układu liniowego i brak wymiernych
wspólnych zer dwóch form na $\mathbb{P}^1\times\mathbb{P}^1$), powtórzonych niezależnie w Sage i sympy.
Tym samym minimalny stopień nad $\Q$ wynosi $6$ (UDOWODNIONE przy cytowanym twierdzeniu, z krokami obliczeniowymi).
Przeszkoda jest czysto arytmetyczna: nad $\Q(\sqrt{94})$ i nad pewnym ciałem sześciennym istnieją
automorfizmy stopnia $4$ realizujące tę bijekcję (wynik obliczeniowy).
Otwarte pozostają m.in.\ ręczny dowód kroków obliczeniowych, dokładne odwołanie bibliograficzne
oraz opis wszystkich rozwiązań stopnia $6$.
